\part{I: Logic as Constraint}\label{part-i-logic-as-constraint}

\hypertarget{chapter-1-the-myth-of-a-single-logic}{%
\chapter*{Chapter 1 --- The Myth of a Single Logic}
\addcontentsline{toc}{chapter}{Chapter 1 --- The Myth of a Single Logic}
\markboth{Chapter 1 --- The Myth of a Single Logic}{Chapter 1 --- The Myth of a Single Logic}\label{chapter-1-the-myth-of-a-single-logic}}

The question this book opens with is not ``which logic is correct.'' It's:

\begin{quote}
What must remain invariant when representations are combined,
transported, refined, contradicted, or extended?
\end{quote}

That question doesn't presuppose an answer about how many truth values
there are, whether contradictions explode, or whether every proposition
is either true or false. It only presupposes that representations get
combined, moved, refined, sometimes contradicted, and sometimes extended
--- and asks what survives those operations. Logic, on this view, is the
study of what a given regime of transformation preserves.

\textbf{Consequence as preservation under admissible transformation.} Let a
\emph{state} be whatever a representational system currently commits to --- a
set of accepted propositions, a configuration of evidence, a model. Write
\(x \leadsto_R y\) for ``state \(y\) is reachable from state \(x\) under the
transformations a regime \(R\) permits.'' Logical consequence, in this
vocabulary, is:
\[\Gamma \models_R \varphi \quad\text{iff}\quad \text{every } R\text{-admissible extension of a state satisfying } \Gamma \text{ also satisfies } \varphi.\]
This is not a new definition invented to be exotic. It's the same shape
as the ordinary one --- consequence as truth-preservation across models ---
with one change: which transformations count as admissible is now a
parameter, \(R\), rather than fixed in advance.

\textbf{Classical logic is a point, not the origin.} Classical consequence is
what you get when \(R\) is set to its most stringent possible value: no
tolerance for a proposition and its negation both holding (no gluts), no
tolerance for a proposition holding neither way (no gaps), and
unrestricted propagation --- once a contradiction is admitted anywhere, every
proposition becomes admissible everywhere. That last clause is
\emph{ex falso quodlibet}, and it is worth noticing that it is a substantive
commitment about \(R\), not a logical necessity. Setting \(R\) to be less
strict --- tolerating gluts, tolerating gaps, restricting how far a local
failure propagates --- doesn't produce a rival to classical logic in the
sense of a competing claim about the same subject matter. It produces a
different regime, and classical logic is the member of that family
where the dial is turned to maximum strictness in every direction at
once. Later chapters (starting in Chapter 3, and given full formal
treatment in Part VI) will make this literal: the space of possible
regimes is a real object, not a metaphor, and classical logic is a
identifiable, describable point within it --- not the space's center, and
not its default.

\textbf{What this buys, and what it doesn't yet.} Reframing consequence this
way doesn't resolve any dispute about which logic to use for a given
purpose. What it does is stop the question from being framed as ``classical
logic, plus some deviant alternatives'' --- a framing that makes every
non-classical system look like a patch or a concession. If classical
logic is one regime among many, non-classical systems aren't deviations
requiring justification; they're other points that need locating, the
same way classical logic does. That relocation is this book's actual
project, and it starts, in the next chapter, from further back than
truth values themselves.

\hypertarget{chapter-2-distinction-before-truth}{%
\chapter*{Chapter 2 --- Distinction Before Truth}
\addcontentsline{toc}{chapter}{Chapter 2 --- Distinction Before Truth}
\markboth{Chapter 2 --- Distinction Before Truth}{Chapter 2 --- Distinction Before Truth}\label{chapter-2-distinction-before-truth}}

The ordinary starting point for logic takes truth values as given:
propositions are things that are true or false (or, in richer systems,
some fixed number of other things), and logic studies how those values
combine. This chapter argues that starting point already presupposes
something that itself needs an account --- and that the account, once
given, changes what a truth value turns out to be.

\textbf{The primitive act is distinction, not valuation.} Before anything can
be true or false, something has to be \emph{distinguished} --- carved out as a
determinate content that a truth value could even apply to. Asking
whether \(p\) is true already presupposes that \(p\) names a specific
distinction: this, as opposed to not-this, within some domain. That
prior act --- drawing the distinction --- is not itself a truth-evaluable
event. It's what makes truth-evaluation possible afterward.

\textbf{What follows from taking this seriously.} If distinction comes first,
a proposition's truth value shouldn't be treated as a primitive label
attached to it from outside. It should be \emph{read off} whatever process
generated the distinction and gathered evidence bearing on it. This is
not a promissory note for later --- it's exactly what Part VI does, in full
formal detail: Chapter 19 reconstructs the four Belnap-Dunn values not as
a stipulated alphabet but as the four possible outcomes of asking two
independent questions --- is there evidence \emph{for} this distinction holding,
and is there evidence \emph{against} it --- with \(N\), \(T\), \(F\), \(B\) arising as a
consequence of the support data rather than as an assumed starting
vocabulary. A glut isn't a fifth, exotic truth value bolted onto a
two-valued base; it's simply what ``both kinds of evidence present''
already means once truth values are treated as outputs rather than
inputs.

\textbf{Why this matters for the rest of the book.} Every non-classical system
covered in this book --- paraconsistent, intuitionistic, many-valued,
relevant, fuzzy --- is usually introduced by first fixing a set of truth
values and then giving rules for how they combine. That's a defensible
way to present a formal system, but it obscures a question worth asking
first: \emph{why these values, and not others?} Once truth values are treated
as consequences of a prior act of evidence-gathering over a distinction,
that question has an answer that doesn't have to be stipulated separately
for every logic in turn --- the values on offer are exactly the values a
given evidential structure can produce, and different logics turn out to
differ in what evidential structure they're built to track, not in an
arbitrary choice of alphabet.

\hypertarget{chapter-3-consistency-and-closure}{%
\chapter*{Chapter 3 --- Consistency and Closure}
\addcontentsline{toc}{chapter}{Chapter 3 --- Consistency and Closure}
\markboth{Chapter 3 --- Consistency and Closure}{Chapter 3 --- Consistency and Closure}\label{chapter-3-consistency-and-closure}}

``Consistent'' is usually treated as a single property a system either has
or lacks. This chapter argues that collapsing several genuinely different
notions into one word is a large part of why paraconsistency looks
paradoxical from the outside, and separates out the notions that actually
do the work.

\textbf{Local consistency.} A single region of a representational system ---
one patch, one context, one moment --- is locally consistent with respect
to a proposition \(p\) if it doesn't itself hold both \(p\) and its negation.
This is a property of one piece of the system in isolation, and nothing
about it says anything yet about how that piece relates to any other.

\textbf{Global consistency.} The whole system, taken together, is globally
consistent if no proposition is jointly supported and refuted once every
region's contributions are combined. This is a strictly different claim
from local consistency at every region individually --- a system can be
locally consistent everywhere and still fail to be globally consistent,
if different regions individually support and refute the same
proposition without either one containing a contradiction on its own.
(Part VI works this distinction out in full: Chapter 21's Glut-Free
Reconstruction Obstruction is precisely the condition under which local
data of this kind admits a globally consistent reading at all.)

\textbf{Constraint closure.} A regime \(R\) (Chapter 1's admissibility relation)
is closed under a class of transformations if performing them never
takes an admissible state outside admissibility. This is what actually
governs whether a local failure spreads: classical explosion is the
claim that \(R\) is closed under ``infer anything from a contradiction'' ---
which is one specific, and maximally permissive, choice of what counts
as an admissible move once a contradiction is present. Restricting
closure --- refusing to admit that move --- is not refusing to reason. It's
choosing a different, more conservative, and equally coherent notion of
what a contradiction is permitted to license.

\textbf{Representational coherence.} Separate again from all three above:
whether the pieces contributing to a proposition's support actually
\emph{agree} with each other, not merely whether they license the same
verdict. Two regions can both support \(p\) being true without their
support being the same support --- Chapter 20's coherence defect gives this
its formal treatment, and it is a genuinely fourth notion, not a
restatement of the first three. A system can be locally consistent
everywhere, globally consistent, closed under a conservative propagation
rule, and still be incoherent in this sense: agreeing on verdicts for the
wrong reasons.

Keeping these four separate --- local consistency, global consistency,
closure, coherence --- is not pedantry. Confusing any two of them is
exactly what makes claims like ``tolerating contradiction is dangerous''
sound like a single, self-evident worry, when in fact the danger (if
there is one) lives specifically in unrestricted closure, and neither
tolerating local gluts nor tolerating global disagreement nor tolerating
incoherent-but-agreeing witnesses is the same danger, or necessarily a
danger at all.

\hypertarget{chapter-4-failure-without-collapse}{%
\chapter*{Chapter 4 --- Failure Without Collapse}
\addcontentsline{toc}{chapter}{Chapter 4 --- Failure Without Collapse}
\markboth{Chapter 4 --- Failure Without Collapse}{Chapter 4 --- Failure Without Collapse}\label{chapter-4-failure-without-collapse}}

This chapter's title is meant as a thesis, not a description: a
representational system can fail --- locally, even repeatedly --- without
that failure collapsing the system into uselessness. What follows is an
account of the three different things ``failure'' can mean, and the
condition that actually separates recoverable failure from catastrophic
failure.

\textbf{Three modes of failure.}

\emph{Contradiction} is a local matter: a proposition and its negation both
supported, somewhere. On its own, per Chapter 3, this says nothing about
the health of the system as a whole.

\emph{Obstruction} is a structural matter, not a semantic one: local pieces of
evidence that individually make sense fail to assemble into a coherent
global picture, whether or not any single piece is itself contradictory.
Chapter 20's coherence defect is the formal name for this; it can occur
with no contradiction present at all.

\emph{Refusal} is neither of the above --- it's a deliberate constitutive act,
the decision that some transformation is not admissible, made in advance
rather than discovered as a failure after the fact. A regime that refuses
a move isn't failing to reason its way past a contradiction; it's
declining, as a matter of what the regime \emph{is}, to treat that move as
available.

\textbf{The distinction that actually matters: inconsistency versus
triviality.} Classical logic's core commitment is that these three modes
of failure, whenever they occur, must be treated identically --- as
signals that something has gone wrong that licenses anything at all to
follow. That identification is the actual source of explosion, and it's
worth stating as its own claim, separable from the rest of classical
logic's machinery: \emph{inconsistency} (a contradiction is present somewhere)
is not the same fact as \emph{triviality} (literally everything is now
derivable). A system can be inconsistent --- even irreducibly, permanently
inconsistent at some point --- without being trivial, provided its closure
conditions (Chapter 3) don't license the move from one to the other.
Priest's paraconsistent logics~\cite{priest2006} are, in this vocabulary, exactly regimes
where that licensing move has been withheld: contradiction is admitted as
a fact a system can contain, while triviality remains exactly as
catastrophic and exactly as avoidable as it always was.

\textbf{What ``failure without collapse'' actually requires.} Not tolerance for
error, and not a weakened notion of truth. It requires a regime whose
closure conditions don't route every local failure into global
collapse --- which is a structural property of \(R\), specifiable and
checkable, not a permissive attitude adopted toward contradiction for its
own sake. Part VI gives this requirement its sharpest form: Chapter 22's
repair machinery treats a glut as something with a repair cost, not as an
alarm demanding total system reset, and Chapter 21's separation of
contradiction from obstruction is precisely the demonstration that
``something has failed'' is not one fact but several, each with its own,
independently manageable consequences.

\hypertarget{status-of-this-draft}{%
\section{Status of this draft}\label{status-of-this-draft}}

Part I is written to be readable on its own --- a reader encountering only
these four chapters gets a complete philosophical case for treating
admissibility, distinction, and constrained closure as prior to truth
values and classical consistency --- while forward-referencing Part VI
honestly rather than pretending its formal results belong here. Nothing
in this draft is a new technical claim; the technical content is Part
VI's, established in Chapters 19--22. What's new here is the exposition:
stating the philosophical case in a form that motivates Part VI's
machinery before the reader has seen it, rather than only in retrospect.
One thing to watch when Parts II--V are drafted: Chapter 3's four-way
consistency/closure/coherence distinction is doing real setup work for
Part II's treatment of non-normal and impossible worlds, and Part V's
identity-under-transport material will likely want to lean on Chapter
2's distinction-before-truth framing rather than restate it --- worth
checking those connections hold once those parts exist, rather than
assuming they will.
